% Part: second-order-logic
% Chapter: metatheory
% Section: compactness

\documentclass[../../../include/open-logic-section]{subfiles}

\begin{document}

\olfileid[fr]{sol}{met}{com}

\olsection{La logique du second ordre n'est pas compacte}

\begin{explain}
Un ensemble d'!!{sentence}s~$\Gamma$ est
\emph{finiment satisfaisable} si chacun de ses
sous-ensembles finis est satisfaisable.
La logique du premier ordre possède la propriété
suivante : tout ensemble finiment satisfaisable
d'!!{sentence}s est satisfaisable. On l'appelle
\emph{compacité}. Une formulation équivalente
concerne la conséquence logique : si
$\Gamma\Entails !A$, il existe déjà un
sous-ensemble fini $\Gamma_0\subseteq\Gamma$
tel que $\Gamma_0\Entails !A$.
Sous cette forme, la compacité découle aussitôt
du théorème de complétude. En effet, de
$\Gamma\Entails !A$ on déduit
$\Gamma\Proves !A$, et une !!{derivation}
ne peut utiliser qu'un nombre fini
d'!!{sentence}s de~$\Gamma$.

La compacité échoue pour la logique du second
ordre sous sa sémantique standard. Certains
ensembles d'!!{sentence}s du second ordre
sont finiment satisfaisables sans être
satisfaisables; ils impliquent alors un
!!{sentence}~$!A$ sans qu'aucun de leurs
sous-ensembles finis ne l'implique.
\end{explain}

\begin{thm}
\ollabel{thm:sol-not-compact}
La logique du second ordre n'est pas compacte.
\end{thm}

\begin{proof}
Rappelons l'!!{sentence}
\[
\fn{Inf} \ident \lexists[u][(\lforall[x][\lforall[y][(\eq[u(x)][u(y)]
      \lif \eq[x][y])]] \land \lexists[y][\lforall[x][\eq/[y][u(x)]]])].
\]
Il est satisfait dans une !!{structure} si et
seulement si son !!{domain} est infini au sens
de Dedekind. Dans le cadre usuel admettant
l'axiome du choix dénombrable, cela équivaut
à l'infinitude ordinaire.

Pour $n\ge 1$, soit $!A^{\ge n}$ un
!!{sentence} qui affirme que le !!{domain}
possède au moins $n$ !!{element}s distincts.
Par exemple, en écrivant une suite de
quantificateurs existentiels, on peut poser
\[
!A^{\ge n}\ident
\lexists[x_1][\dots\lexists[x_n][
\bigwedge_{1\le i<j\le n}\eq/[x_i][x_j]]].
\]
Considérons l'ensemble d'!!{sentence}s
\[
\Gamma=\{\lnot\fn{Inf},
!A^{\ge1},!A^{\ge2},!A^{\ge3},\dots\}.
\]
Il est finiment satisfaisable. En effet, pour
un sous-ensemble fini~$\Gamma_0$, prenons
$k\ge1$ au moins aussi grand que tous les
indices~$n$ tels que $!A^{\ge n}\in\Gamma_0$.
Une !!{structure} dont le !!{domain} a exactement
$k$ !!{element}s satisfait~$\Gamma_0$.

En revanche, $\Gamma$ n'est pas satisfaisable.
Si $\Sat{M}{\lnot\fn{Inf}}$, le !!{domain}
est fini dans le cadre précisé, disons de
cardinal~$k$. Mais alors
$\Sat/{M}{!A^{\ge k+1}}$, tandis que cet
!!{sentence} appartient à~$\Gamma$.
\end{proof}

\begin{editorial}
La source étiquette ce théorème par le même
identifiant que celui d'indécidabilité; nous
employons un identifiant distinct.
Son exemple de $!A^{\ge n}$ abrège par des
points de suspension les inégalités entre
!!{variable}s. Nous rendons explicite la
conjonction sur toutes les paires $i<j$, afin
que les $n$ valeurs soient bien distinctes.
Dans l'argument de
satisfaisabilité finie, la borne des indices
concerne $\Gamma_0$, et non l'ensemble
infini~$\Gamma$ comme l'écrit la source.
Enfin, l'équivalence entre l'infinitude au
sens de Dedekind et l'infinitude ordinaire
utilise ici le choix dénombrable, comme dans
la section précédente.
\end{editorial}

\begin{prob}
Donnez un ensemble~$\Gamma$ et un
!!{sentence}~$!A$ tels que
$\Gamma\Entails !A$, mais que
$\Gamma_0\Entails/ !A$ pour tout
sous-ensemble fini $\Gamma_0\subseteq\Gamma$.
\end{prob}

\end{document}
